% Por qué las curvas de indiferencia (a) tienen pendiente negativa, (b) dan más utilidad
% cuanto más alejadas del origen, (c) no se cortan y (d) son convexas, con el axioma
% que da cada propiedad. Curvas del tipo x1·x2 = k.
% Autor: Víctor Gutiérrez Marcos
\documentclass[tikz,border=4pt]{standalone}
% ==========================================================================
% tikz-tcee.tex — Estilos comunes de los gráficos TikZ del temario
% Autor: Víctor Gutiérrez Marcos
%
% Cada gráfico es un fichero standalone en fuentes/<tema>/graficos/:
%
%   \documentclass[tikz,border=4pt]{standalone}
%   % ==========================================================================
% tikz-tcee.tex — Estilos comunes de los gráficos TikZ del temario
% Autor: Víctor Gutiérrez Marcos
%
% Cada gráfico es un fichero standalone en fuentes/<tema>/graficos/:
%
%   \documentclass[tikz,border=4pt]{standalone}
%   % ==========================================================================
% tikz-tcee.tex — Estilos comunes de los gráficos TikZ del temario
% Autor: Víctor Gutiérrez Marcos
%
% Cada gráfico es un fichero standalone en fuentes/<tema>/graficos/:
%
%   \documentclass[tikz,border=4pt]{standalone}
%   \input{../../../latex/tikz-tcee}
%   \begin{document}
%   \begin{tikzpicture}[tcee]
%     \ejes{Q}{P}{6}{5}
%     \draw[curva1] (0.5,4.5) -- (5,0.8) node[etiqueta,right] {$D$};
%     \capa{2}{\draw[curva2] ...;}   % en el vídeo aparece en el paso 2
%   \end{tikzpicture}
%   \end{document}
%
% Paleta: la de la web (morado, dorado, salvia) más tres colores de apoyo
% distinguibles en blanco y negro por el trazo.
% ==========================================================================
\usepackage[T1]{fontenc}
\usepackage{newpxtext,newpxmath}
\usepackage{xcolor}
\usepackage{tikz,pgfplots}
\pgfplotsset{compat=1.18}
\usetikzlibrary{arrows.meta,calc,positioning,patterns,decorations.pathreplacing,intersections,shapes.misc,backgrounds}

\definecolor{tceemorado}{HTML}{5F2987}
\definecolor{tceedorado}{HTML}{B8860B}
\definecolor{tceeverde}{HTML}{3C7D3A}
\definecolor{tceeazul}{HTML}{1F5F99}
\definecolor{tceerojo}{HTML}{B22222}
\definecolor{tceegris}{HTML}{767171}
\definecolor{tceesalvia}{HTML}{E2EFD9}

% Capas para el vídeo: \capa{n}{…} dibuja su contenido solo a partir del paso n.
% En el PDF, la web y Word se ve todo; generar-video.py compila el gráfico una
% vez por paso con \def\tceepaso{k} para que cada elemento aparezca al narrarlo.
\providecommand{\tceepaso}{1000}
\newcommand{\capa}[2]{\ifnum#1>\tceepaso\relax\else#2\fi}

\tikzset{
  tcee/.style={line cap=round,line join=round,font=\small,>={Stealth[length=2.2mm]}},
  eje/.style={->,thick,black},
  curva1/.style={very thick,tceemorado},
  curva2/.style={very thick,tceeazul},
  curva3/.style={very thick,tceeverde},
  curva4/.style={very thick,tceedorado},
  curvanueva/.style={very thick,tceerojo},
  desplazada/.style={very thick,dashed},
  guia/.style={thin,dashed,tceegris},
  punto/.style={circle,fill=black,inner sep=1.3pt},
  etiqueta/.style={font=\small,inner sep=2pt},
  flecha/.style={->,thick,tceerojo},
  area/.style={fill=tceemorado,fill opacity=0.18,draw=none},
  area2/.style={fill=tceedorado,fill opacity=0.22,draw=none},
}

% \ejes{etiqueta x}{etiqueta y}{largo x}{alto y}
\newcommand{\ejes}[4]{%
  \draw[eje] (0,0) -- (#3,0) node[below left=2pt and 0pt] {#1};
  \draw[eje] (0,0) -- (0,#4) node[left] {#2};
  \node[below left] at (0,0) {$0$};}
% \proyeccion{(x,y)}{etq x}{etq y}: líneas guía a los ejes con etiquetas
\newcommand{\proyeccion}[3]{%
  \draw[guia] let \p1=#1 in (\x1,0) node[below] {#2} |- (0,\y1) node[left] {#3};}

%   \begin{document}
%   \begin{tikzpicture}[tcee]
%     \ejes{Q}{P}{6}{5}
%     \draw[curva1] (0.5,4.5) -- (5,0.8) node[etiqueta,right] {$D$};
%     \capa{2}{\draw[curva2] ...;}   % en el vídeo aparece en el paso 2
%   \end{tikzpicture}
%   \end{document}
%
% Paleta: la de la web (morado, dorado, salvia) más tres colores de apoyo
% distinguibles en blanco y negro por el trazo.
% ==========================================================================
\usepackage[T1]{fontenc}
\usepackage{newpxtext,newpxmath}
\usepackage{xcolor}
\usepackage{tikz,pgfplots}
\pgfplotsset{compat=1.18}
\usetikzlibrary{arrows.meta,calc,positioning,patterns,decorations.pathreplacing,intersections,shapes.misc,backgrounds}

\definecolor{tceemorado}{HTML}{5F2987}
\definecolor{tceedorado}{HTML}{B8860B}
\definecolor{tceeverde}{HTML}{3C7D3A}
\definecolor{tceeazul}{HTML}{1F5F99}
\definecolor{tceerojo}{HTML}{B22222}
\definecolor{tceegris}{HTML}{767171}
\definecolor{tceesalvia}{HTML}{E2EFD9}

% Capas para el vídeo: \capa{n}{…} dibuja su contenido solo a partir del paso n.
% En el PDF, la web y Word se ve todo; generar-video.py compila el gráfico una
% vez por paso con \def\tceepaso{k} para que cada elemento aparezca al narrarlo.
\providecommand{\tceepaso}{1000}
\newcommand{\capa}[2]{\ifnum#1>\tceepaso\relax\else#2\fi}

\tikzset{
  tcee/.style={line cap=round,line join=round,font=\small,>={Stealth[length=2.2mm]}},
  eje/.style={->,thick,black},
  curva1/.style={very thick,tceemorado},
  curva2/.style={very thick,tceeazul},
  curva3/.style={very thick,tceeverde},
  curva4/.style={very thick,tceedorado},
  curvanueva/.style={very thick,tceerojo},
  desplazada/.style={very thick,dashed},
  guia/.style={thin,dashed,tceegris},
  punto/.style={circle,fill=black,inner sep=1.3pt},
  etiqueta/.style={font=\small,inner sep=2pt},
  flecha/.style={->,thick,tceerojo},
  area/.style={fill=tceemorado,fill opacity=0.18,draw=none},
  area2/.style={fill=tceedorado,fill opacity=0.22,draw=none},
}

% \ejes{etiqueta x}{etiqueta y}{largo x}{alto y}
\newcommand{\ejes}[4]{%
  \draw[eje] (0,0) -- (#3,0) node[below left=2pt and 0pt] {#1};
  \draw[eje] (0,0) -- (0,#4) node[left] {#2};
  \node[below left] at (0,0) {$0$};}
% \proyeccion{(x,y)}{etq x}{etq y}: líneas guía a los ejes con etiquetas
\newcommand{\proyeccion}[3]{%
  \draw[guia] let \p1=#1 in (\x1,0) node[below] {#2} |- (0,\y1) node[left] {#3};}

%   \begin{document}
%   \begin{tikzpicture}[tcee]
%     \ejes{Q}{P}{6}{5}
%     \draw[curva1] (0.5,4.5) -- (5,0.8) node[etiqueta,right] {$D$};
%     \capa{2}{\draw[curva2] ...;}   % en el vídeo aparece en el paso 2
%   \end{tikzpicture}
%   \end{document}
%
% Paleta: la de la web (morado, dorado, salvia) más tres colores de apoyo
% distinguibles en blanco y negro por el trazo.
% ==========================================================================
\usepackage[T1]{fontenc}
\usepackage{newpxtext,newpxmath}
\usepackage{xcolor}
\usepackage{tikz,pgfplots}
\pgfplotsset{compat=1.18}
\usetikzlibrary{arrows.meta,calc,positioning,patterns,decorations.pathreplacing,intersections,shapes.misc,backgrounds}

\definecolor{tceemorado}{HTML}{5F2987}
\definecolor{tceedorado}{HTML}{B8860B}
\definecolor{tceeverde}{HTML}{3C7D3A}
\definecolor{tceeazul}{HTML}{1F5F99}
\definecolor{tceerojo}{HTML}{B22222}
\definecolor{tceegris}{HTML}{767171}
\definecolor{tceesalvia}{HTML}{E2EFD9}

% Capas para el vídeo: \capa{n}{…} dibuja su contenido solo a partir del paso n.
% En el PDF, la web y Word se ve todo; generar-video.py compila el gráfico una
% vez por paso con \def\tceepaso{k} para que cada elemento aparezca al narrarlo.
\providecommand{\tceepaso}{1000}
\newcommand{\capa}[2]{\ifnum#1>\tceepaso\relax\else#2\fi}

\tikzset{
  tcee/.style={line cap=round,line join=round,font=\small,>={Stealth[length=2.2mm]}},
  eje/.style={->,thick,black},
  curva1/.style={very thick,tceemorado},
  curva2/.style={very thick,tceeazul},
  curva3/.style={very thick,tceeverde},
  curva4/.style={very thick,tceedorado},
  curvanueva/.style={very thick,tceerojo},
  desplazada/.style={very thick,dashed},
  guia/.style={thin,dashed,tceegris},
  punto/.style={circle,fill=black,inner sep=1.3pt},
  etiqueta/.style={font=\small,inner sep=2pt},
  flecha/.style={->,thick,tceerojo},
  area/.style={fill=tceemorado,fill opacity=0.18,draw=none},
  area2/.style={fill=tceedorado,fill opacity=0.22,draw=none},
}

% \ejes{etiqueta x}{etiqueta y}{largo x}{alto y}
\newcommand{\ejes}[4]{%
  \draw[eje] (0,0) -- (#3,0) node[below left=2pt and 0pt] {#1};
  \draw[eje] (0,0) -- (0,#4) node[left] {#2};
  \node[below left] at (0,0) {$0$};}
% \proyeccion{(x,y)}{etq x}{etq y}: líneas guía a los ejes con etiquetas
\newcommand{\proyeccion}[3]{%
  \draw[guia] let \p1=#1 in (\x1,0) node[below] {#2} |- (0,\y1) node[left] {#3};}

\begin{document}
\begin{tikzpicture}[tcee,
    titulo/.style={font=\small\bfseries,align=center,anchor=base},
    pie/.style={font=\footnotesize,align=center,anchor=north}]
  \newcommand{\ejesp}{\ejes{$x_1$}{$x_2$}{4.4}{4.2}}
  % ================= (a) pendiente negativa: monotonía =================
  \begin{scope}[shift={(0,0)}]
    \capa{1}{\node[titulo] at (2.2,4.6) {(a) Pendiente negativa (monotonía)};}
    \capa{1}{\ejesp}
    \capa{2}{\node[punto] (a0) at (2,2) {};}
    \capa{2}{\draw[guia] (0,2) -- (4.2,2);}
    \capa{2}{\draw[guia] (2,0) -- (2,4);}
    \capa{3}{\fill[area] (2,2) rectangle (4.2,4);}
    \capa{3}{\node[etiqueta,align=center] at (3.1,3.3) {más de\\ todo: $\succ x^0$};}
    \capa{4}{\fill[area2] (0,0) rectangle (2,2);}
    \capa{4}{\node[etiqueta,align=center] at (1,0.75) {menos de\\ todo: $\prec x^0$};}
    \capa{5}{\draw[curva1] plot[domain=1.0:4.0,samples=50] (\x,{4/\x}) node[etiqueta,above] {$I(x^0)$};}
    \capa{5}{\node[punto] at (2,2) {};}
    \capa{5}{\node[etiqueta,above right] at (2,2) {$x^0$};}
    \capa{5}{\node[pie] at (2.2,-0.5) {la curva solo puede pasar por los cuadrantes\\ noroeste y sudeste de $x^0$};}
  \end{scope}
  % ================= (b) más alejadas, más utilidad: monotonía =================
  \begin{scope}[shift={(6.8,0)}]
    \capa{6}{\node[titulo] at (2.2,4.6) {(b) Más alejadas, más utilidad (monotonía)};}
    \capa{6}{\ejesp}
    \capa{7}{\draw[curva1] plot[domain=0.4:4.0,samples=50] (\x,{1.5/\x}) node[etiqueta,above right] {$U_0$};}
    \capa{7}{\draw[curva2] plot[domain=0.75:4.0,samples=50] (\x,{3/\x}) node[etiqueta,above right] {$U_1$};}
    \capa{7}{\draw[curva3] plot[domain=1.25:4.0,samples=50] (\x,{5/\x}) node[etiqueta,above right] {$U_2$};}
    \capa{8}{\draw[guia] (0,0) -- (3.4,3.4);}
    \capa{8}{\node[punto] (bA) at (1.2247,1.2247) {};}
    \capa{8}{\node[etiqueta,below right] at (bA) {$A$};}
    \capa{8}{\node[punto] (bB) at (1.7321,1.7321) {};}
    \capa{8}{\node[etiqueta,below right] at (bB) {$B$};}
    \capa{8}{\node[punto] (bC) at (2.2361,2.2361) {};}
    \capa{8}{\node[etiqueta,below right] at (bC) {$C$};}
    \capa{9}{\node[pie] at (2.2,-0.5) {$C$ tiene más de ambos bienes que $B$, y $B$ que $A$:\\ $C\succ B\succ A$, luego $U_2>U_1>U_0$};}
  \end{scope}
  % ================= (c) no se cortan: transitividad =================
  \begin{scope}[shift={(0,-7.0)}]
    \capa{10}{\node[titulo] at (2.2,4.6) {(c) No se cortan (transitividad)};}
    \capa{10}{\ejesp}
    \capa{11}{\draw[curva1] plot[domain=1.0:4.0,samples=50] (\x,{4/\x}) node[etiqueta,below] {$U_1$};}
    \capa{12}{\draw[curva2] plot[domain=0.6:4.0,samples=50] (\x,{1.2+1.6/\x}) node[etiqueta,above] {$U_2$};}
    \capa{12}{\node[punto] (cA) at (2,2) {};}
    \capa{12}{\node[etiqueta,above right] at (cA) {$A$};}
    \capa{13}{\node[punto] (cB) at (3,1.7333) {};}
    \capa{13}{\node[etiqueta,above] at (cB) {$B$};}
    \capa{13}{\node[punto] (cC) at (3,1.3333) {};}
    \capa{13}{\node[etiqueta,below] at (cC) {$C$};}
    \capa{14}{\node[pie] at (2.2,-0.5) {$B\sim A$ y $A\sim C\Rightarrow B\sim C$ (transitividad),\\ pero $B$ tiene más bien 2 que $C$,\\ luego $B\succ C$: contradicción};}
  \end{scope}
  % ================= (d) convexas: convexidad =================
  \begin{scope}[shift={(6.8,-7.0)}]
    \capa{15}{\node[titulo] at (2.2,4.6) {(d) Convexas (convexidad)};}
    \capa{15}{\ejesp}
    \capa{16}{\fill[area] plot[domain=0.976:4.0,samples=50] (\x,{4/\x}) -- (4.0,4.1) -- (0.976,4.1) -- cycle;}
    \capa{16}{\draw[curva1] plot[domain=0.976:4.0,samples=50] (\x,{4/\x}) node[etiqueta,above right] {$I(x')$};}
    \capa{17}{\node[punto] (dx1) at (1.2,3.3333) {};}
    \capa{17}{\node[etiqueta,left] at (dx1) {$x'$};}
    \capa{17}{\node[punto] (dx2) at (3.2,1.25) {};}
    \capa{17}{\node[etiqueta,below left] at (dx2) {$x''$};}
    \capa{18}{\draw[thick,dashed,tceerojo] (dx1) -- (dx2);}
    \capa{18}{\node[punto,fill=tceerojo] (dm) at (2.2,2.2917) {};}
    \capa{18}{\node[etiqueta,above right,text=tceerojo] at (dm) {$\lambda x'+(1-\lambda)x''$};}
    \capa{19}{\node[pie] at (2.2,-0.5) {las combinaciones de dos cestas indiferentes son\\ al menos tan buenas: el conjunto $B(x')$ es convexo};}
  \end{scope}
\end{tikzpicture}
\end{document}
